\documentclass[a4paper,12pt]{article}
\usepackage[T2A]{fontenc}
\usepackage[utf8]{inputenc}
\usepackage[russian]{babel}
\usepackage{amsmath,amssymb,mathtools}
\usepackage{DejaVuSans}
\renewcommand{\familydefault}{\sfdefault}
\usepackage{icomma}
\usepackage[a4paper,margin=20mm,headheight=25pt,headsep=8mm,footskip=12mm]{geometry}
\usepackage{microtype}
\usepackage{needspace}
\usepackage{tikz}
\usetikzlibrary{calc,arrows.meta,positioning,shapes.geometric}
\usepackage{tabularx,booktabs,array}
\usepackage{enumitem}
\usepackage{fancyhdr}
\usepackage{xcolor}
\usepackage{setspace}
\usepackage{hyperref}
\usepackage{parskip}

\definecolor{accent}{HTML}{2C6E6F}
\definecolor{darkaccent}{HTML}{214F50}
\definecolor{textgray}{HTML}{292F30}
\definecolor{softgray}{HTML}{687475}
\definecolor{linegray}{HTML}{CCD9D8}
\hypersetup{colorlinks=true,linkcolor=darkaccent,urlcolor=darkaccent}
\setlength{\parindent}{0pt}
\setstretch{1.17}
\setlist[itemize]{leftmargin=1.5em,itemsep=0.25em,topsep=0.28em}
\setlist[enumerate]{leftmargin=1.7em,itemsep=0.85em,topsep=0.45em}
\pagestyle{fancy}
\fancyhf{}
\fancyhead[L]{\small\color{softgray} Смешанные числа и математические модели}
\fancyhead[R]{\small\color{softgray} 13.10.26}
\fancyfoot[C]{\small\color{softgray}\thepage}
\renewcommand{\headrulewidth}{0.3pt}
\renewcommand{\headrule}{\hbox to\headwidth{\color{linegray}\leaders\hrule height \headrulewidth\hfill}}

\newcommand{\mixed}[3]{#1\,\frac{#2}{#3}}
\newcommand{\sectiontitle}[1]{%
  \Needspace{5\baselineskip}%
  \vspace{0.45em}{\large\bfseries\color{darkaccent}#1}\par
  \vspace{0.18em}{\color{linegray}\hrule}\vspace{0.48em}}
\newcommand{\keyidea}[1]{\textbf{Главная мысль.} #1}
\newcommand{\note}[1]{\textbf{Важно.} #1}
\newcommand{\ex}[1]{\textbf{Пример.} #1}

\begin{document}
\color{textgray}
\setlength{\abovedisplayskip}{8pt}
\setlength{\belowdisplayskip}{8pt}
\setlength{\abovedisplayshortskip}{5pt}
\setlength{\belowdisplayshortskip}{5pt}

\pagenumbering{roman}
\begin{titlepage}
\thispagestyle{empty}
\centering
\vspace*{27mm}
{\small\color{softgray} МАТЕМАТИКА}\par
\vspace{9mm}
{\Huge\bfseries\color{darkaccent} Чек-лист}\par
\vspace{8mm}
{\LARGE\bfseries Смешанные числа и текстовые задачи}\par
\vspace{3mm}
{\large вычисления, уравнения, движение и совместная работа}\par
\vspace{17mm}
{\large 13.10.26}\par
\vspace{10mm}
{\normalsize Тренировочный блок: Путь А · 10 заданий}\par
\vfill
{\large\bfseries Лёвин Артём Александрович}\par
\vspace{2mm}
{\normalsize эксклюзивно для Нади Клименко}\par
\vspace{19mm}
\begin{tikzpicture}[remember picture,overlay]
\draw[accent,line width=1.05pt]
  ($(current page.south west)+(18mm,20mm)$)
  .. controls ($(current page.south west)+(65mm,30mm)$)
  and ($(current page.south east)+(-66mm,10mm)$)
  .. ($(current page.south east)+(-18mm,20mm)$);
\end{tikzpicture}
\end{titlepage}
\pagenumbering{arabic}

\sectiontitle{1. Смешанные числа: сначала упростить, затем вычислить}
\keyidea{При сложении и вычитании смешанных чисел отдельно учитывайте целые и дробные части. Для дробных частей сначала найдите общий знаменатель.}
\[
\mixed{2}{3}{10}+\mixed{1}{1}{5}
=3+\frac3{10}+\frac2{10}
=\boxed{\mixed{3}{1}{2}}.
\]
Если дробной части уменьшаемого недостаточно, возьмите одну единицу из целой части:
\[
\begin{aligned}
\mixed{4}{1}{6}-\mixed{2}{3}{4}
&=\mixed{4}{2}{12}-\mixed{2}{9}{12}\\
&=\mixed{3}{14}{12}-\mixed{2}{9}{12}
=\boxed{\mixed{1}{5}{12}}.
\end{aligned}
\]
\note{Одна занятая единица превращается в \(\frac{12}{12}\), поэтому \(\mixed{4}{2}{12}=\mixed{3}{14}{12}\).}

\sectiontitle{2. Целое число, десятичная дробь и скобки}
Целое число можно представить обыкновенной дробью с нужным знаменателем. Это особенно полезно перед умножением и делением дробей:
\[
1-\frac4{17}=\frac{17}{17}-\frac4{17}=\boxed{\frac{13}{17}},
\qquad 1=\frac{17}{17}.
\]
Десятичные дроби удобно переводить в обыкновенные с последующим сокращением:
\[
0{,}35=\frac{35}{100}=\frac7{20},\qquad
2{,}4=\mixed{2}{4}{10}=\mixed{2}{2}{5}.
\]
\ex{Сначала вычислите выражение в скобках:}
\[
\begin{aligned}
\mixed{3}{1}{6}-\Bigl(\mixed{1}{3}{4}-\mixed{1}{1}{3}\Bigr)
&=\mixed{3}{1}{6}-\left(\frac9{12}-\frac4{12}\right)\\
&=\mixed{3}{2}{12}-\frac5{12}
=\boxed{\mixed{2}{3}{4}}.
\end{aligned}
\]
\note{Если несколько действий содержат дроби, сначала выполните самые удобные действия и сократите дроби. Не ищите общий знаменатель для всего выражения заранее.}

\sectiontitle{3. Как перевести фразу в уравнение}
\keyidea{Обозначьте неизвестное число через \(x\) и сначала запишите словесную связь. Для нахождения неизвестного слагаемого используется вычитание, для неизвестного уменьшаемого --- сложение.}

\begin{tabularx}{\linewidth}{@{}X>{\raggedright\arraybackslash}p{0.36\linewidth}@{}}
\toprule
Текст & Уравнение \\
\midrule
К числу прибавили \(\mixed{1}{1}{2}\), получили \(\mixed{3}{1}{4}\). & \(x+\mixed{1}{1}{2}=\mixed{3}{1}{4}\)\\[2pt]
Число уменьшили на \(\mixed{1}{8}{9}\), получили \(\mixed{6}{2}{3}\). & \(x-\mixed{1}{8}{9}=\mixed{6}{2}{3}\)\\
\bottomrule
\end{tabularx}
\[
 x=\mixed{3}{1}{4}-\mixed{1}{1}{2}=\boxed{\mixed{1}{3}{4}},
\qquad
 x=\mixed{6}{2}{3}+\mixed{1}{8}{9}=\boxed{\mixed{8}{5}{9}}.
\]
\note{После вычисления подставьте найденное число в исходное уравнение. Слова «увеличили» и «уменьшили» относятся к действию над исходным числом.}

\sectiontitle{4. Текстовая задача: сначала математическая модель}
\keyidea{Прежде чем считать, запишите, какие величины относятся к каждому этапу. Словам «на столько-то меньше» соответствует вычитание, а словам «всего за два этапа» --- сложение.}

\ex{За первый месяц ребёнок прибавил \(\mixed{1}{3}{25}\)~кг --- это на \(\frac7{20}\)~кг больше, чем за второй месяц. Сколько килограммов он прибавил за оба месяца?}
\[
\begin{aligned}
\text{Второй месяц:}\quad
&\mixed{1}{3}{25}-\frac7{20}
=\mixed{1}{12}{100}-\frac{35}{100}=\frac{77}{100}\ \text{кг};\\
\text{Всего:}\quad
&\mixed{1}{3}{25}+\frac{77}{100}
=\mixed{1}{12}{100}+\frac{77}{100}
=\boxed{\mixed{1}{89}{100}\ \text{кг}}.
\end{aligned}
\]
\note{Проверка смысла: прирост за второй месяц меньше прироста за первый на \(\frac7{20}\)~кг.}

\sectiontitle{5. Движение по течению и против течения}
Для равномерного движения \(v=\frac{S}{t}\). Скорость лодки \emph{в стоячей воде} --- её собственная скорость \(v_0\):
\[
 v_{\text{по}}=v_0+u,\qquad v_{\text{против}}=v_0-u,
\]
где \(u\) --- скорость течения, \(v_0>u\).

\ex{Лодка проходит 3750~м за 9~мин в стоячей воде. Скорость течения \(\mixed{20}{5}{6}\)~м/мин.}
\[
 v_0=\frac{3750}{9}=\frac{1250}{3}=\mixed{416}{2}{3}\ \text{м/мин}.
\]
\[
\begin{aligned}
 v_{\text{по}}&=\mixed{416}{4}{6}+\mixed{20}{5}{6}
=\boxed{\mixed{437}{1}{2}\ \text{м/мин}},\\
 v_{\text{против}}&=\mixed{416}{4}{6}-\mixed{20}{5}{6}
=\boxed{\mixed{395}{5}{6}\ \text{м/мин}}.
\end{aligned}
\]
\note{Единицы времени должны совпадать. При движении по течению скорость больше собственной, против течения --- меньше.}

\sectiontitle{6. Три этапа и известные суммы соседних этапов}
Обозначьте результаты трёх рейсов \(a,b,c\). Если известно \(a+b\), \(b+c\) и \(a+b+c\), удобнее сначала найти \emph{крайние} рейсы:
\[
 c=(a+b+c)-(a+b),\qquad
 a=(a+b+c)-(b+c),\qquad
 b=(a+b)-a.
\]
\ex{Всего перевезли 20~т; за два первых рейса --- \(\mixed{14}{5}{12}\)~т, за два последних --- \(\mixed{13}{11}{20}\)~т.}
\[
\begin{aligned}
 c&=20-\mixed{14}{5}{12}=\mixed{5}{7}{12}\ \text{т},\\
 a&=20-\mixed{13}{11}{20}=\mixed{6}{9}{20}\ \text{т},\\
 b&=\mixed{14}{5}{12}-\mixed{6}{9}{20}=\mixed{7}{29}{30}\ \text{т}.
\end{aligned}
\]
\note{Проверьте одновременно три равенства: \(a+b\), \(b+c\) и \(a+b+c\). Так легче заметить арифметическую ошибку.}

\sectiontitle{7. Совместная работа: производительность}
\keyidea{Если вся работа равна 1, а время её выполнения --- \(t\), то производительность равна \(p=\frac1t\). При совместной работе производительности складываются.}

\ex{Игорь и Паша красят забор за 3~ч; Паша и Володя --- за 6~ч; Володя и Игорь --- за 4~ч. Сколько минут потребуется всем троим?}

Пусть \(i,p,v\) --- производительности мальчиков (частей забора в час). Тогда
\[
 i+p=\frac13,\qquad p+v=\frac16,\qquad v+i=\frac14.
\]
При сложении трёх равенств производительность \emph{каждого} учитывается дважды:
\[
 2(i+p+v)=\frac13+\frac16+\frac14=\frac34,
\qquad i+p+v=\frac38.
\]
\[
 t=\frac1{3/8}=\frac83\ \text{ч}=\boxed{160\ \text{мин}}.
\]
\note{Если время дано для пар, складывание парных производительностей удваивает общую производительность. Перед нахождением времени обязательно разделите сумму на 2.}

\clearpage
\thispagestyle{plain}
\sectiontitle{Алгоритм: от текстового условия к вычислению}
\begin{center}
\begin{tikzpicture}[
  node distance=12mm and 15mm,
  >={Latex[length=2.4mm]},
  line/.style={draw=darkaccent,line width=0.8pt,-Latex},
  start/.style={ellipse,draw=accent,line width=0.9pt,align=center,minimum width=44mm,minimum height=12mm},
  process/.style={rectangle,rounded corners=2mm,draw=accent,line width=0.8pt,align=center,text width=97mm,minimum height=14mm,inner sep=3mm},
  cond/.style={diamond,aspect=2.8,draw=accent,line width=0.8pt,align=center,text width=41mm,inner sep=1mm},
  fork/.style={rectangle,rounded corners=2mm,draw=accent,line width=0.8pt,align=center,text width=51mm,minimum height=18mm,inner sep=3mm},
  every node/.style={font=\small}
]
\node[start] (begin) {Прочитайте условие};
\node[process,below=of begin] (data) {Обозначьте искомые величины.\\ Подпишите единицы измерения и известные данные.};
\node[cond,below=15mm of data] (type) {Какая связь описана?};
\node[fork,below left=17mm and 8mm of type] (sum) {Сумма, разность,\\ «больше / меньше»\\ \(\Rightarrow\) выражение или уравнение};
\node[fork,below right=17mm and 8mm of type] (rate) {Движение, работа\\ \(v=S/t\), \(p=A/t\)\\ \(\Rightarrow\) свяжите величины};
\node[process,below=47mm of type] (model) {Запишите математическую модель.\\ Вычислите, соблюдая порядок действий.};
\node[process,below=of model] (check) {Проверьте ответ: единицы, знак, разумность результата, подстановка.};
\node[start,below=of check] (end) {Запишите ответ};
\draw[line] (begin)--(data);
\draw[line] (data)--(type);
\draw[line] (type.south west) -| (sum.north);
\draw[line] (type.south east) -| (rate.north);
\coordinate (join) at ([yshift=8mm]model.north);
\draw[draw=darkaccent,line width=0.8pt] (sum.south) |- (join);
\draw[draw=darkaccent,line width=0.8pt] (rate.south) |- (join);
\draw[line] (join) -- (model.north);
\draw[line] (model)--(check);
\draw[line] (check)--(end);
\end{tikzpicture}
\end{center}

\clearpage
\sectiontitle{Тренировочный блок · Путь А}
\textbf{10 заданий.} Выполняйте по порядку: сначала вычисления и уравнения, затем текстовые задачи. Записывайте математическую модель и ответ. Решения в этом разделе намеренно отсутствуют.

\begin{enumerate}
\item Вычислите: \(\displaystyle 1-\frac5{18}\).
\item Вычислите: \(\displaystyle \mixed{4}{7}{12}+\mixed{2}{5}{18}\).
\item Вычислите: \(\displaystyle \mixed{6}{1}{8}-\mixed{2}{5}{12}\).
\item Представьте десятичные дроби в виде обыкновенных и вычислите:
\[2{,}4-\Bigl(\mixed{1}{3}{10}+0{,}35\Bigr).\]
\item К числу прибавили \(\mixed{2}{1}{6}\) и получили \(\mixed{6}{3}{4}\). Составьте уравнение и найдите число.
\item Число уменьшили на \(\mixed{1}{7}{15}\) и получили \(\mixed{3}{2}{5}\). Составьте уравнение и найдите число.
\item Лодка в стоячей воде проходит 1500~м за 4~мин. Скорость течения равна \(\mixed{12}{1}{2}\)~м/мин. Найдите скорости лодки по течению и против течения в м/мин.
\item За первый месяц ребёнок прибавил \(\mixed{1}{7}{20}\)~кг, а за второй --- на \(\frac3{10}\)~кг меньше. Сколько килограммов он прибавил за два месяца?
\item Самосвал перевёз за три рейса 18~т груза. За первые два рейса он перевёз \(\mixed{11}{3}{5}\)~т, за последние два --- \(\mixed{12}{1}{4}\)~т. Сколько тонн пришлось на каждый рейс отдельно?
\item \textbf{Повышенная сложность.} Игорь и Паша вместе выполняют работу за 4~ч; Паша и Володя --- за 4~ч; Игорь и Володя --- за 6~ч. За сколько минут выполнят ту же работу все трое вместе, если работают с постоянной производительностью?
\end{enumerate}

\clearpage
\sectiontitle{Ответы к тренировочному блоку}
\begin{center}
\renewcommand{\arraystretch}{1.65}
\begin{tabularx}{\linewidth}{@{}>{\centering\arraybackslash}p{18mm}X@{}}
\toprule
Номер & Ответ\\
\midrule
1 & \(\displaystyle \frac{13}{18}\)\\
2 & \(\displaystyle \mixed{6}{31}{36}\)\\
3 & \(\displaystyle \mixed{3}{17}{24}\)\\
4 & \(\displaystyle \frac34\)\\
5 & \(\displaystyle x=\mixed{4}{7}{12}\)\\
6 & \(\displaystyle x=\mixed{4}{13}{15}\)\\
7 & По течению: \(\mixed{387}{1}{2}\)~м/мин; против течения: \(\mixed{362}{1}{2}\)~м/мин.\\
8 & \(\displaystyle \mixed{2}{2}{5}\)~кг\\
9 & Первый: \(\mixed{5}{3}{4}\)~т; второй: \(\mixed{5}{17}{20}\)~т; третий: \(\mixed{6}{2}{5}\)~т.\\
10 & \(180\)~мин\\
\bottomrule
\end{tabularx}
\end{center}
\vspace{7mm}
\textbf{Для самопроверки текстовых задач:} проверьте, что сумма результатов за три рейса равна общему грузу; скорость по течению больше собственной скорости; в задаче на совместную работу каждая пара участников учитывается ровно один раз в своём условии.
\vspace{3mm}
\sectiontitle{Самопроверка по материалу занятия}
\begin{itemize}
\item[\(\square\)] Я умею приводить дроби к общему знаменателю и занимать единицу.
\item[\(\square\)] Я перевожу десятичные дроби в обыкновенные и учитываю скобки.
\item[\(\square\)] Я составляю уравнение из словесного условия.
\item[\(\square\)] Я могу назвать смысл каждого действия в текстовой задаче.
\item[\(\square\)] Я различаю собственную скорость и скорость течения.
\item[\(\square\)] Я нахожу результаты трёх этапов по соседним суммам.
\item[\(\square\)] Я складываю производительности, учитывая повторный подсчёт участников.
\end{itemize}


\end{document}
